\chapter{Собираем процессор}
\label{ch:cpu}

Мы изучили сумматоры, регистры, счётчики, память и конечные автоматы. Все
вместе они образуют процессор. В этой главе мы соберём его целиком:
восьмиразрядный, с десятью командами, который выполняет программу из ПЗУ и
умеет, например, вычислять числа Фибоначчи. Он занимает несколько десятков
строк Verilog и работает в Logra. Это завершение книги: каждый узел, который мы
изучили, понадобится здесь.

\section{Из чего состоит процессор}

Любой процессор делится на две части: \term{тракт данных} (datapath), где
хранятся и преобразуются числа, и \term{устройство управления}, которое
решает, что делать на каждом такте.

\begin{figure}[H]
\centering
\begin{tikzpicture}[font=\sffamily\small, bx/.style={draw, minimum width=24mm, minimum height=11mm, fill=white, rounded corners=2pt, align=center}]
  \node[bx, fill=FpgaLight] (pc) at (0,3) {Счётчик\\команд PC};
  \node[bx] (rom) at (4,3) {ПЗУ\\программы};
  \node[bx] (ir) at (8,3) {Регистр\\команды IR};
  \node[bx, fill=green!10] (ctl) at (8,0.8) {Автомат\\управления};
  \node[bx, fill=FpgaLight] (acc) at (0,0.8) {Аккумулятор\\ACC};
  \node[bx] (alu) at (0,-1.5) {Сумматор /\\вычитатель};
  \node[bx] (ram) at (4,-1.5) {ОЗУ\\данных 16$\times$8};
  \draw[-{Stealth}] (pc) -- node[above] {адрес} (rom);
  \draw[-{Stealth}] (rom) -- node[above] {команда} (ir);
  \draw[-{Stealth}] (ir) -- (ctl);
  \draw[-{Stealth}] (ctl.west) -- ++(-1.0,0) |- (acc.east) node[pos=0.25, above, font=\sffamily\footnotesize] {управление};
  \draw[-{Stealth}] (acc) -- (alu);
  \draw[-{Stealth}] (ram) -- (alu);
  \draw[-{Stealth}] (alu.west) -- ++(-0.7,0) |- (acc.west);
  \draw[-{Stealth}] (acc.south east) -- (ram.north west);
\end{tikzpicture}
\caption{Структура процессора. Слева --- тракт данных, справа --- устройство управления}
\end{figure}

\begin{itemize}
  \item \textbf{Счётчик команд} (PC) хранит адрес следующей команды и после каждой
        выборки увеличивается на единицу --- это обычный счётчик из
        главы~\ref{ch:regs}.
  \item \textbf{ПЗУ программы} хранит коды команд. Для процессора оно --- таблица:
        по адресу выдаётся команда (глава~\ref{ch:sync}).
  \item \textbf{Регистр команды} (IR) запоминает команду на время выполнения.
  \item \textbf{Аккумулятор} (ACC) --- главный регистр: в нём лежит результат
        каждой операции.
  \item \textbf{Сумматор/вычитатель} --- АЛУ из главы~\ref{ch:arith}.
  \item \textbf{ОЗУ данных} --- 16 ячеек по 8 бит для промежуточных значений.
  \item \textbf{Автомат управления} --- конечный автомат с двумя состояниями:
        «выбрать команду» и «выполнить», который вырабатывает все управляющие
        сигналы.
\end{itemize}

Такую организацию, где память программ и память данных разделены, называют
\term{гарвардской}. В настольных компьютерах память общая (архитектура фон
Неймана), но для ПЛИС раздельные блоки проще: у каждого свои порты.

\begin{diybox}[Найдите блоки процессора в RTL]
\textbf{Задание.} Сопоставьте блоки рисунка структуры и элементы схемы, которую строит синтезатор.

\textbf{В Logra.} Карточка «На основе примера» --- «Маленький процессор: fetch и execute», либо \code{tiny\_cpu.v} из книги; вкладка «Код» $\to$ «Синтезировать»; вкладка RTL (группа «Проверить»).

\textbf{Ожидаемый результат.} На схеме видны триггеры \code{pc}, \code{ir}, \code{acc}, \code{out}, память \code{ram}, сумматор и мультиплексоры. Схема велика (около пятидесяти элементов), поэтому ищите нужное по именам, а не по положению.

\textbf{Как проверить.} Сверьте названия с рис.~\ref{fig:n-cpu-rtl}: каждое имя в рамке должно найтись на RTL.

\stars{1} повторите: найдите \code{pc}, \code{ir}, \code{acc}. \stars{2} измените: переименуйте \code{acc} в \code{a} и посмотрите, что изменилось на схеме. \stars{3} проект: этап «Охранная панель» --- процессор как сервисный блок; критерий приёмки --- составлена таблица «функция панели $\to$ блок процессора, который её выполнит».
\end{diybox}

\section{Система команд}

Команда занимает один байт: старшие 4 бита --- код операции, младшие 4 ---
операнд (число или адрес ячейки ОЗУ).

\begin{table}[H]
\centering
\caption{Система команд процессора}
\small
\begin{tabular}{clll}
\toprule
\textbf{Код} & \textbf{Мнемоника} & \textbf{Действие} & \textbf{Операнд} \\
\midrule
0 & NOP & ничего не делать & --- \\
1 & LDI $k$ & ACC $\leftarrow k$ & число 0--15 \\
2 & ADD $a$ & ACC $\leftarrow$ ACC + RAM[$a$] & адрес \\
3 & STA $a$ & RAM[$a$] $\leftarrow$ ACC & адрес \\
4 & LDA $a$ & ACC $\leftarrow$ RAM[$a$] & адрес \\
5 & SUB $a$ & ACC $\leftarrow$ ACC $-$ RAM[$a$] & адрес \\
6 & JMP $a$ & PC $\leftarrow a$ & адрес команды \\
7 & JZ $a$ & если ACC = 0, то PC $\leftarrow a$ & адрес команды \\
8 & OUT & показать ACC на выходе & --- \\
9 & HLT & остановиться & --- \\
\bottomrule
\end{tabular}
\end{table}

Короткие имена команд --- \term{мнемоники}. Программу на языке мнемоник
называют ассемблерным текстом, а их перевод в числа выполняет ассемблер. У
нас он занимает минуту ручной работы: LDI 1 --- это \code{0x11}, STA 1 ---
\code{0x31} и так далее.

\begin{diybox}[Свой ассемблер вручную]
\textbf{Задание.} Составьте программу, которая складывает 1 и 2 и выводит результат, и переведите её в коды.

\textbf{В Logra.} Вкладка «Код»: в функции \code{rom} замените ветки \code{case} значениями \code{11 30 12 20 80 90}. Это \code{LDI 1; STA 0; LDI 2; ADD 0; OUT; HLT}. Вкладка «Диаграммы»: \code{pc}, \code{acc}, \code{out}.

\textbf{Ожидаемый результат.} После \code{OUT} значение \code{out} равно 3, затем процессор останавливается (\code{HLT}).

\textbf{Как проверить.} Диаграммы: число циклов 20; \code{pc} идёт $0,1,2,3,4,5$ и больше не растёт.

\stars{1} повторите: программа $1+2$. \stars{2} измените: вычислите $5-3$ командами \code{SUB}. \stars{3} проект: этап «Охранная панель» --- программа с командами процессора; критерий приёмки --- составлена таблица адресов и кодов, по ней вручную предсказано значение \code{out}, и оно совпало с диаграммой.
\end{diybox}

\section{Программа: числа Фибоначчи}

Каждое число Фибоначчи равно сумме двух предыдущих: $0, 1, 1, 2, 3, 5, 8,
13, \ldots$. Хранить будем два числа $a$ и $b$ в ячейках ОЗУ 0 и 1 и временное
значение $t$ в ячейке 2. На каждом шаге показываем $a$, затем
вычисляем $a+b$ и сдвигаем пару: $a \leftarrow b$, $b \leftarrow t$.

\begin{table}[H]
\centering
\caption{Программа вычисления чисел Фибоначчи}
\small
\begin{tabular}{rllll}
\toprule
\textbf{Адрес} & \textbf{Код} & \textbf{Команда} & \textbf{Смысл} \\
\midrule
0 & \code{0x11} & LDI 1 & \\
1 & \code{0x31} & STA 1 & $b = 1$ \\
2 & \code{0x10} & LDI 0 & \\
3 & \code{0x30} & STA 0 & $a = 0$ \\
4 & \code{0x40} & LDA 0 & \textbf{начало цикла} \\
5 & \code{0x80} & OUT & показать $a$ \\
6 & \code{0x21} & ADD 1 & ACC $= a + b$ \\
7 & \code{0x32} & STA 2 & $t = a + b$ \\
8 & \code{0x41} & LDA 1 & \\
9 & \code{0x30} & STA 0 & $a = b$ \\
10 & \code{0x42} & LDA 2 & \\
11 & \code{0x31} & STA 1 & $b = t$ \\
12 & \code{0x64} & JMP 4 & в начало цикла \\
\bottomrule
\end{tabular}
\end{table}

\begin{diybox}[Числа Фибоначчи по шагам]
\textbf{Задание.} Выполните программу из таблицы на диаграмме и проверьте каждое число.

\textbf{В Logra.} Проект с \code{tiny\_cpu.v}; вкладка «Диаграммы»: число циклов 140, сигналы \code{pc}, \code{ir}, \code{acc}, \code{out}.

\textbf{Ожидаемый результат.} В \code{out} по очереди появляются $0, 1, 1, 2, 3, 5, 8, 13, \ldots$; один проход цикла занимает 18 тактов.

\textbf{Как проверить.} Диаграммы: на рис.~\ref{fig:n-cpu-wave} видно, как \code{pc} возвращается к \code{0x4} по команде \code{JMP 4}.

\stars{1} повторите: первые шесть чисел. \stars{2} измените: начните с $a=1$, $b=2$ (измените константы в первых четырёх командах) и предскажите ряд. \stars{3} проект: этап «Охранная панель» --- счётчик попыток ввода кода; критерий приёмки --- программа выводит 3, 2, 1, 0 и останавливается.
\end{diybox}

\section{Описание на Verilog}

Процессор записан в одном модуле. В комментариях сначала перечислены все
команды; ПЗУ описано функцией с оператором \code{case} --- это таблица кодов
из предыдущего раздела.

\vfile{\codepath/09_cpu/tiny_cpu.v}{code/09\_cpu/tiny\_cpu.v}

Рассмотрим устройство управления. Регистр \code{phase} хранит текущее
состояние автомата. В фазе 0 процессор читает команду из ПЗУ в регистр
\code{ir} и увеличивает счётчик \code{pc}, в фазе 1 --- выполняет её: оператор
\code{case} по коду операции выбирает, что записать в аккумулятор, в ОЗУ или
в счётчик команд. Поэтому каждая команда занимает два такта. На рис.~\ref{fig:n-cpu-rtl} те же блоки
показаны с теми именами, что и в коде: по нему удобно читать вкладку RTL, где
автоматическая раскладка показывает полсотни элементов и сотни связей.

\begin{figure}[H]
\centering
\begin{tikzpicture}[font=\sffamily\footnotesize,
  bx/.style={draw=FpgaBlue, thick, rounded corners=2pt, minimum width=24mm, minimum height=11mm, align=center, fill=white},
  reg/.style={bx, fill=FpgaLight},
  mem/.style={bx, fill=orange!12, draw=FpgaOrange},
  ctlb/.style={bx, fill=green!10, draw=FpgaGreen, minimum height=9mm},
  d/.style={-{Stealth[length=2mm]}, thick, FpgaDark},
  c/.style={-{Stealth[length=2mm]}, thick, FpgaGreen!60!black, dashed},
  t/.style={font=\sffamily\footnotesize, text=FpgaDark, inner sep=1.5pt}]
  % верхний ряд: выборка команды
  \node[reg] (pc)  at (0,0)    {Счётчик\\команд\\{\ttfamily pc[3:0]}};
  \node[mem] (rom) at (4.4,0)  {ПЗУ\\программы\\{\ttfamily rom(pc)}};
  \node[reg] (ir)  at (8.8,0)  {Регистр\\команды\\{\ttfamily ir[7:0]}};
  \draw[d] (pc) -- node[t, above] {адрес} (rom);
  \draw[d] (rom) -- node[t, above] {команда} (ir);
  % устройство управления
  \node[ctlb, minimum width=100mm] (ctl) at (4.4,-2.6) {Автомат управления: {\ttfamily phase} (выборка / исполнение), {\ttfamily case (op)}};
  \draw[d] (ir.south) -- node[t, right] {{\ttfamily op = ir[7:4]}} (ir.south |- ctl.north);
  \draw[d] (ctl.west) -- ++(-0.6,0) |- node[t, pos=0.25, left] {{\ttfamily pc <= arg} (JMP, JZ)} (pc.west);
  % нижний ряд: тракт данных
  \node[reg] (acc) at (0,-5.6)  {Аккумулятор\\{\ttfamily acc[7:0]}};
  \node[bx]  (alu) at (4.4,-5.6) {Сумматор /\\вычитатель};
  \node[mem] (ram) at (8.8,-5.6) {ОЗУ данных\\{\ttfamily ram[0:15]}};
  \draw[c] (acc.north |- ctl.south) -- node[t, right] {загрузка} (acc.north);
  \draw[c] (alu.north |- ctl.south) -- node[t, right] {$+$ или $-$} (alu.north);
  \draw[c] (ram.north |- ctl.south) -- node[t, right] {запись} (ram.north);
  \draw[d] ([yshift=2mm]acc.east) -- node[t, above] {acc} ([yshift=2mm]alu.west);
  \draw[d] ([yshift=-2mm]alu.west) -- node[t, below] {итог} ([yshift=-2mm]acc.east);
  \draw[d] (ram.west) -- node[t, above] {ram[arg]} (alu.east);
  % адрес ОЗУ из младших бит команды
  \draw[d] (ir.east) -- ++(1.0,0) |- node[t, pos=0.25, right] {{\ttfamily arg}} (ram.east);
  % запись аккумулятора в ОЗУ (STA) и вывод (OUT)
  \draw[d] (acc.south) -- ++(0,-0.9) -| node[t, pos=0.25, below] {acc $\to$ ram (STA)} (ram.south);
  \draw[d] (acc.west) -- ++(-1.0,0) node[t, left] {{\ttfamily out[7:0]} (OUT)};
\end{tikzpicture}
\caption{Структура процессора с именами из кода. Сплошные стрелки --- данные, пунктир --- управляющие сигналы}
\label{fig:n-cpu-rtl}
\end{figure}


\begin{diybox}[Автомат управления: фаза и команда]
\textbf{Задание.} Разберите, как регистр \code{phase} разделяет работу на выборку и исполнение.

\textbf{В Logra.} \code{tiny\_cpu.v}; вкладка «Диаграммы»: сигналы \code{phase}, \code{pc}, \code{ir}, \code{acc}.

\textbf{Ожидаемый результат.} \code{phase} чередуется $0,1,0,1,\ldots$; \code{pc} растёт только в фазе 0; \code{acc} меняется только в фазе 1.

\textbf{Как проверить.} Диаграммы: убедитесь, что любая команда занимает ровно два такта. На RTL найдите регистр \code{phase}.

\stars{1} повторите: найдите такт, на котором \code{ir} получает \code{0x31}. \stars{2} измените: сделайте \code{HLT}-состояние видимым --- выведите \code{halted} как порт. \stars{3} проект: этап «Охранная панель» --- управление процессором; критерий приёмки --- по сигналу \code{rst} процессор возвращается в состояние выборки с \code{pc = 0}.
\end{diybox}

\section{Запуск в Logra}

Откройте вкладку «Диаграммы», в настройках сценария поставьте число циклов
140 и запустите симуляцию. Выберите в списке сигналов \code{pc}, \code{ir},
\code{acc} и \code{out}.

\shot{n-cpu-wave}{Процессор выполняет программу: счётчик \code{pc} идёт по адресам, в \code{ir} видны коды команд, в \code{out} появляются числа 0, 1, 2, 3, 5, 8}{fig:n-cpu-wave}

На диаграмме (рис.~\ref{fig:n-cpu-wave}) можно проследить за работой процессора:
счётчик \code{pc} последовательно принимает значения 0, 1, 2 и так далее,
регистр \code{ir} содержит коды команд из таблицы (\code{0x11}, \code{0x31},
\code{0x10}, \ldots), в аккумуляторе \code{acc} появляются значения, а в
\code{out} после каждой команды \code{OUT} появляется очередное число
Фибоначчи. Поскольку первые два числа (1 и 1) совпадают, на диаграмме выход
остаётся равным 1 дольше. После адреса 12 команда \code{JMP 4} возвращает
счётчик в начало цикла: поэтому \code{pc} на диаграмме опять показывает
\code{0x4}.

\begin{ideabox}
Процессор --- всего лишь конечный автомат, у которого состояние содержит
в себе регистры (PC, IR, ACC), а память работает как ПЗУ программы. Всё,
что нужно, мы уже изучили: счётчик, сумматор, регистр, дешифратор команд,
мультиплексор, память и автомат. Современные процессоры отличаются
масштабом (конвейер, кэш, сотни команд), но не принципом.
\end{ideabox}

\begin{diybox}[Замедлить процессор для светодиодов]
\textbf{Задание.} Выведите \code{out} на светодиоды платы так, чтобы числа было видно глазом.

\textbf{В Logra.} Вкладка «Код»: добавьте делитель с импульсом разрешения \code{tick} (глава~\ref{ch:clocks}) и выполняйте шаг процессора только при \code{tick = 1}. Подключите \code{out[5:0]} к \code{led}. Вкладка «Выводы»: назначьте светодиоды.

\textbf{Ожидаемый результат.} Светодиоды показывают ряд Фибоначчи в двоичном виде, смена раз в 0,5~с.

\textbf{Как проверить.} Диаграммы: при уменьшенном делителе \code{out} меняется каждые $N$ тактов. Затем вкладка «Лаборатория» или плата.

\stars{1} повторите: значения \code{out} на диаграмме. \stars{2} измените: скорость вдвое быстрее. \stars{3} проект: этап «Охранная панель» --- вывод на плату; критерий приёмки --- светодиоды меняются не чаще раза в 0,5~с, все триггеры тактируются одним \code{clk}.
\end{diybox}

\section{Куда развивать дальше}

Эту модель легко расширять, и каждое расширение --- хорошее задание.

\begin{itemize}
  \item \textbf{Условные переходы.} Добавьте флаг «перенос» и команду \code{JC},
        чтобы программа могла остановиться, когда число перестало
        умещаться в 8 бит.
  \item \textbf{Подпрограммы и стек.} Команды \code{CALL} и \code{RET} с
        небольшой стековой памятью.
  \item \textbf{Ввод-вывод.} Подключите \code{out} к светодиодам платы, а
        \code{in} --- к кнопкам; для видимости замедлите процессор делителем
        частоты (с сигналом разрешения, глава~\ref{ch:clocks}).
  \item \textbf{Конвейер.} Совместите выборку следующей команды с исполнением
        текущей: скорость почти удвоится.
  \item \textbf{Загрузка программы по UART.} Принимайте байты программы по
        UART из главы~\ref{ch:uart} и записывайте их в ПЗУ, сделав его ОЗУ.
\end{itemize}

\begin{diybox}[Команда JNZ]
\textbf{Задание.} Добавьте в процессор команду \code{JNZ a} (код 10 $=$ \code{0xA}): если \code{acc} не равен нулю, то \code{pc <= a}.

\textbf{В Logra.} В блок \code{case (op)} добавьте \code{4'd10: if (acc != 8'd0) pc <= arg;}. В функции \code{rom} запишите программу \code{15 30 11 31 40 80 51 30 A4 90}, то есть \code{LDI 5; STA 0; LDI 1; STA 1; LDA 0; OUT; SUB 1; STA 0; JNZ 4; HLT}.

\textbf{Ожидаемый результат.} На \code{out} по очереди появляются 5, 4, 3, 2, 1, после чего процессор останавливается.

\textbf{Как проверить.} Диаграммы: число циклов 100; \code{out} принимает значения 5 до 1 и больше не меняется; \code{pc} возвращается на 4 по команде \code{JNZ}.

\stars{1} повторите: запустите программу с \code{JNZ}. \stars{2} найдите ошибку: замените \code{JNZ 4} на \code{JNZ 2} и объясните, почему вывод зациклился. \stars{3} проект: этап «Охранная панель» --- блокировка после трёх неверных кодов; критерий приёмки --- программа считает попытки 3, 2, 1, 0 и останавливается.
\end{diybox}

\section{Вопросы и задачи}
\label{sec:tasks-cpu}

Ответы собраны в приложении~\ref{app:answers}.

\begin{taskbox}[Задачи к главе]
\begin{enumerate}
  \item \stars{1}~Сколько тактов нужно, чтобы выполнить один проход цикла программы
        Фибоначчи (адреса 4--12)?
  \item \stars{2}~Составьте программу, которая считает от 5 до 0 и на каждом шаге
        выдаёт число на выход. Подсказка: используйте \code{SUB} и \code{JZ}.
  \item \stars{1}~Сколько ячеек памяти программ можно адресовать 4-разрядным счётчиком
        команд? Что придётся изменить, чтобы увеличить программу вдвое?
  \item \stars{2}~Какое число Фибоначчи окажется последним из тех, что умещаются в 8 бит?
        Что появится на выходе следом?
  \item \stars{1}~Чем отличаются гарвардская архитектура и архитектура фон Неймана?
  \item \stars{2}~Как изменить процессор, чтобы команда выполнялась за один такт?
  \item \stars{2}~Найдите ошибку. В программе из задачи~2 последняя команда цикла записана как
        \code{JMP 2} вместо \code{JMP 3}. Процессор бесконечно выводит число 5. Объясните, почему.
  \item \stars{3}~Проект. Добавьте команду \code{JNZ} и напишите программу, считающую от 5 до 1.
        Критерий приёмки: на диаграмме \code{out} принимает значения 5, 4, 3, 2, 1, затем процессор останавливается.
  \item \stars{3}~Проект «Охранная панель», этап «блокировка». Программа считает неверные попытки ввода
        кода (3, 2, 1, 0) и останавливается при нуле. Критерий приёмки: на диаграмме после третьей
        попытки \code{pc} перестаёт меняться, а \code{out} равно 0.
\end{enumerate}
\end{taskbox}
